% ============================================================================
%  SECCIÓN 5.5: IMPLEMENTACIÓN DEL CRUD PERSISTENTE
% ============================================================================

\section{Implementaci\'on del CRUD persistente}

En este avance se completa la implementaci\'on de las operaciones CRUD
(Create, Read, Update, Delete) con persistencia en Supabase. A diferencia
del avance anterior, donde la operaci\'on de actualizaci\'on estaba
pendiente, ahora se implementan las cinco operaciones b\'asicas
para las entidades principales del sistema.

\subsection{CRUD de escenarios}

La entidad \textit{scenarios} es la m\'as completa del sistema,
implementando todas las operaciones CRUD a trav\'es del hook
\textit{useManageScenarios}:

\subsubsection{Read: consulta de escenarios}

La operaci\'on de lectura se implementa en dos variantes: una para
el cat\'alogo p\'ublico y otra para la gesti\'on administrativa.

\begin{lstlisting}
// Cat\'alogo: solo escenarios activos
async function fetchScenarios():
  Promise<ScenarioWithDetails[]> {
  const { data, error } = await supabase
    .from('scenarios')
    .select(`
      *,
      scenario_images(url, is_primary,
        storage_path, display_order),
      scenario_sports(sports(nombre))
    `)
    .eq('estado', 'activo')
    .order('nombre');

  if (error) throw new Error(error.message);
  return (data ?? []).map(resolveScenarioImages);
}

// Gesti\'on: todos los escenarios (admin/gestor)
async function fetchAllScenarios():
  Promise<ScenarioWithImages[]> {
  const { data, error } = await supabase
    .from('scenarios')
    .select(`
      *,
      scenario_images(url, is_primary,
        storage_path, display_order)
    `)
    .order('nombre');

  if (error) throw new Error(error.message);
  return (data ?? []).map(resolveScenarioImages);
}
\end{lstlisting}

\subsubsection{Create/Update: upsert de escenario}

La operaci\'on de creaci\'on y actualizaci\'on se unifica mediante
\textit{upsert}, que inserta un nuevo registro si no existe o
actualiza el existente:

\begin{lstlisting}
async function upsertScenario(
  scenario: UpsertScenarioData
): Promise<Scenario> {
  const { data, error } = await supabase
    .from('scenarios')
    .upsert({
      ...scenario,
      updated_at: new Date().toISOString(),
    })
    .select()
    .single();

  if (error) throw new Error(error.message);
  return data as Scenario;
}
\end{lstlisting}

\subsubsection{Delete: eliminaci\'on de escenario}

La eliminaci\'on se implementa en dos variantes: una suave
(\textit{soft delete}) que cambia el estado a inactivo, y una
completa (\textit{hard delete}) que elimina el registro y sus
im\'agenes del almacenamiento:

\begin{lstlisting}
// Soft delete: cambia estado a inactivo
async function softDeleteScenario(
  id: string
): Promise<void> {
  const { error } = await supabase
    .from('scenarios')
    .update({
      estado: 'inactivo',
      updated_at: new Date().toISOString(),
    })
    .eq('id', id);

  if (error) throw new Error(error.message);
}

// Hard delete: elimina registro e im\'agenes
async function hardDeleteScenario(
  id: string
): Promise<void> {
  // Eliminar im\'agenes del storage primero
  const { data: images } = await supabase
    .from('scenario_images')
    .select('storage_path')
    .eq('scenario_id', id);

  if (images && images.length > 0) {
    await supabase.storage
      .from('scenario-images')
      .remove(
        images.map((img) => img.storage_path)
      );
  }

  const { error } = await supabase
    .from('scenarios')
    .delete()
    .eq('id', id);

  if (error) throw new Error(error.message);
}
\end{lstlisting}

\subsection{CRUD de eventos}

La entidad \textit{events} implementa las operaciones Create, Read
y Delete. La actualizaci\'on se gestiona mediante la interfaz de
edici\'on del escenario:

\begin{lstlisting}
// Create: crear evento
async function createEvent(
  eventData: CreateEventData
): Promise<Event> {
  const { data, error } = await supabase
    .from('events')
    .insert({
      scenario_id: eventData.scenario_id,
      nombre: eventData.nombre,
      fecha: eventData.fecha,
      hora: eventData.hora || '00:00',
      descripcion: eventData.descripcion || '',
    })
    .select()
    .single();

  if (error) throw new Error(error.message);
  return data as Event;
}

// Read: eventos pr\'oximos
async function fetchUpcomingEvents():
  Promise<EventWithScenario[]> {
  const today = new Date()
    .toISOString().split('T')[0];

  const { data, error } = await supabase
    .from('events')
    .select(`*, scenarios(*)`)
    .gte('fecha', today)
    .order('fecha', { ascending: true })
    .order('hora', { ascending: true })
    .limit(5);

  if (error) throw new Error(error.message);
  return data as EventWithScenario[];
}

// Delete: eliminar evento
async function deleteEvent(
  eventId: string
): Promise<void> {
  const { error } = await supabase
    .from('events')
    .delete()
    .eq('id', eventId);

  if (error) throw new Error(error.message);
}
\end{lstlisting}

\subsection{CRUD de favoritos}

La funcionalidad de favoritos implementa las operaciones Read,
Create y Delete. La actualizaci\'on no es necesaria porque un
favorito solo se agrega o elimina:

\begin{lstlisting}
// Read: listar favoritos del usuario
async function fetchFavorites(
  userId: string
): Promise<FavoriteWithScenario[]> {
  const { data, error } = await supabase
    .from('favorites')
    .select(`
      *,
      scenarios(
        *,
        scenario_images(url, is_primary,
          storage_path)
      )
    `)
    .eq('user_id', userId)
    .order('created_at', { ascending: false });

  if (error) throw new Error(error.message);
  return (data as FavoriteWithScenario[])
    .map((fav) => ({
      ...fav,
      scenarios: resolveScenarioImages(
        fav.scenarios
      ),
    }));
}

// Create: a~nadir favorito
async function addFavorite(
  userId: string, scenarioId: string
): Promise<void> {
  const { error } = await supabase
    .from('favorites')
    .insert({
      user_id: userId,
      scenario_id: scenarioId,
    });

  if (error) throw new Error(error.message);
}

// Delete: eliminar favorito
async function removeFavorite(
  userId: string, scenarioId: string
): Promise<void> {
  const { error } = await supabase
    .from('favorites')
    .delete()
    .eq('user_id', userId)
    .eq('scenario_id', scenarioId);

  if (error) throw new Error(error.message);
}
\end{lstlisting}

\subsection{Sincronizaci\'on del cach\'e}

Despu\'es de cada mutaci\'on (Create, Update o Delete), se
invalidan las consultas de React Query asociadas para forzar la
recarga de datos en la interfaz:

\begin{lstlisting}
// Tras upsert de escenario
queryClient.invalidateQueries({
  queryKey: ['all-scenarios']
});
queryClient.invalidateQueries({
  queryKey: ['scenarios']
});

// Tras agregar/eliminar favorito
queryClient.invalidateQueries({
  queryKey: ['favorites', userId]
});
queryClient.invalidateQueries({
  queryKey: ['favorite', userId, scenarioId]
});

// Tras crear/eliminar evento
queryClient.invalidateQueries({
  queryKey: ['scenario', variables.scenario_id]
});
queryClient.invalidateQueries({
  queryKey: ['upcoming-events']
});
\end{lstlisting}

\subsection{Resumen de operaciones CRUD}

\begin{table}[H]
  \small
  \centering
  \caption{Operaciones CRUD implementadas en el avance 5}
  \contenidoConFuente{%
    \begin{tablaAPA}
      \begin{tabular}{@{}p{1.2cm}p{2.0cm}p{4.6cm}p{5.2cm}p{1.8cm}@{}}
        \toprule
        \textbf{Op.} & \textbf{Entidad} & \textbf{Funci\'on} &
        \textbf{Hook} & \textbf{Estado} \\
        \midrule
        Read & Escenarios & {\footnotesize\ttfamily fetchScenarios} &
        {\footnotesize\ttfamily useScenarios} & OK \\
        \addlinespace
        Read & Escenarios & {\footnotesize\ttfamily fetchAllScenarios} &
        {\footnotesize\ttfamily useAllScenarios} & OK \\
        \addlinespace
        Read & Eventos & {\footnotesize\ttfamily fetchUpcomingEvents} &
        {\footnotesize\ttfamily useUpcomingEvents} & OK \\
        \addlinespace
        Read & Favoritos & {\footnotesize\ttfamily fetchFavorites} &
        {\footnotesize\ttfamily useFavorites} & OK \\
        \addlinespace
        Create & Escenario & {\footnotesize\ttfamily upsertScenario} &
        {\footnotesize\ttfamily useUpsertScenario} & OK \\
        \addlinespace
        Create & Evento & {\footnotesize\ttfamily createEvent} &
        {\footnotesize\ttfamily useCreateEvent} & OK \\
        \addlinespace
        Create & Favorito & {\footnotesize\ttfamily addFavorite} &
        {\footnotesize\ttfamily useToggleFavorite} & OK \\
        \addlinespace
        Update & Escenario & {\footnotesize\ttfamily upsertScenario} &
        {\footnotesize\ttfamily useUpsertScenario} & OK \\
        \addlinespace
        Update & Escenario & {\footnotesize\ttfamily toggleScenarioStatus} &
        {\footnotesize\ttfamily useToggleScenarioStatus} & OK \\
        \addlinespace
        Delete & Escenario & {\footnotesize\ttfamily softDeleteScenario} &
        {\footnotesize\ttfamily useDeleteScenario} & OK \\
        \addlinespace
        Delete & Escenario & {\footnotesize\ttfamily hardDeleteScenario} &
        {\footnotesize\ttfamily useHardDeleteScenario} & OK \\
        \addlinespace
        Delete & Evento & {\footnotesize\ttfamily deleteEvent} &
        {\footnotesize\ttfamily useDeleteEvent} & OK \\
        \addlinespace
        Delete & Favorito & {\footnotesize\ttfamily removeFavorite} &
        {\footnotesize\ttfamily useToggleFavorite} & OK \\
        \bottomrule
      \end{tabular}
    \end{tablaAPA}%
  }{Elaboraci\'on propia en base al c\'odigo fuente del proyecto.}
\end{table}
